\section{Geometric tracking and initial metrics}\label{sec:geometric}

\subsection{Verification of the phase hypotheses}
Choose the graph of Proposition~\ref{prop:graph} on $B_{4r}(V)$.
Set $v=u-\Psi(p,x)$ and define
\begin{equation}\label{eq:geometric-perturbations}
 \begin{split}
 h_g(p,x,v)&=\beta(p,x,\Psi(p,x)+v)-\beta_0(p,x),\\
 h_x(p,x,v)&=C_xh_g(p,x,v),\\
 \F_p(x,v)&=\N_p(x,\Psi+v)-\N_p(x,\Psi)
                         -D_x\Psi(p,x)C_xh_g(p,x,v).
 \end{split}
\end{equation}
The last subtraction follows from the invariance equation
\eqref{eq:invariance}. In these variables the exact normal equation is
$\dot v=-A_pv+\F_p(x,v)$. Its remainder contains no second
spatial derivative of $v$. All three functions in
\eqref{eq:geometric-perturbations} vanish at $v=0$.
Smoothness and \eqref{eq:vanishing-jets} imply, on smaller
coefficient and normal balls,
\begin{equation}\label{eq:geometric-smallness}
 \|D_vh\|+\|D_v\F_p\|\le\epsilon,\qquad
 \|D_xh\|+\|D_x\F_p\|\le K_r\|v\|_{E_\alpha}.
\end{equation}
Normal remainders are measured from $E_\alpha$ to $E_0$.
Uniform bounds through order two follow from H\"older
multiplication and composition estimates on these balls;
compactness of an infinite-dimensional normal ball is not used.

For phase selection we need a complete finite extension.
Choose a $K$-invariant smooth coefficient cutoff $\chi_r$
equal to one on $B_{2r}$ and supported in $B_{3r}$, and put
\begin{equation}\label{eq:cutoff}
 \widehat\beta_0=b_p+\chi_r(\beta_0-b_p),\qquad
 \widehat h_g=\chi_rh_g,\qquad
 \widehat h_x=C_x\widehat h_g,\qquad
 \widehat\F_p=\chi_r\F_p.
\end{equation}
The supported differences are extended by zero.
The central drift outside the support is retained.
To see the relevant bounds explicitly, let $B(x)=\beta_0(p,x)-b_p$.
Proposition~\ref{prop:moment} gives, on the support,
\[
 |B|\le Cr^2,\qquad |DB|\le Cr,\qquad |D^2B|\le C.
\]
Since $|D^j\chi_r|\le C_jr^{-j}$, the same bounds hold
for $\widehat\beta_0-b_p$ through order two.
For $\widehat f(x)=C_x\widehat\beta_0(x)$,
$|D\widehat f|\le C(|p|+r^2)$.
The full right-trivialized group derivative also contains
$D\widehat\beta_0$, and hence admits the bound
\begin{equation}\label{eq:finite-L}
 L\le C(|p|+r).
\end{equation}
The bracket term in \eqref{eq:finite-linearization} obeys
the same bound. Outside the support the finite field is
the complete central linear drift, and the group coefficient
is bounded everywhere.

Multiplication by $\chi_r$ preserves the small normal
derivative bounds in \eqref{eq:geometric-smallness}.
Coefficient derivatives of the cutoff introduce powers
of $r^{-1}$, but their factors vanish at $v=0$, so the
first coefficient derivatives remain bounded by
$K_r\|v\|$. Second derivative constants are finite for
the chosen radius. Proposition~\ref{prop:phase} does not
require these constants to remain bounded as $r\to0$.
There is no dependence on a finite orbit rank or spectral
ratio. We can therefore choose the parameter and coefficient
radii to make $L$ small relative to $\sigma$, then the
normal radius to make $\epsilon$ small, and finally the
initial displacement threshold $\delta$.

\subsection{Proof of the main theorem}
\begin{proof}[Proof of Theorem~\ref{thm:main}]
Propositions~\ref{prop:graph} and \ref{prop:moment} give
part~(i). Equations~\eqref{eq:semigroup},
\eqref{eq:geometric-perturbations} and
\eqref{eq:cutoff} verify Proposition~\ref{prop:phase}
with $\alpha=7/8$, uniformly in $p$.
It supplies a global solution of the extended system,
normal exponential decay and a group-valued phase.

Before any exit from the geometric domain, uniqueness
identifies this solution with the sliced HYM system.
Confinement prevents coefficient exit. The quadratic
graph bound and
$\|v(s)\|_{h^{17/8}}\le CD e^{-cs}$ prevent normal exit
after reducing $r$ and $\delta$. The group coefficient
is bounded, so $g$ and $g^{-1}$ cannot escape in finite
time. The mild continuation argument with the integrable
smoothing kernel excludes a remaining finite-time exit
in the initial H\"older space. Thus the geometric
coordinate interval is infinite.

The phase estimate gives $|x(s)-x_*(s)|\le CD e^{-cs}$.
Require $C\delta<r/2$. Since $|x(s)|\le r$, the
selected trajectory remains inside $B_{2r}$, where
the finite extension agrees with the exact graph field.
The orbit-preserving conclusion of Proposition~\ref{prop:phase}
gives $x_*(s)=g_*(s)x_b$. Lemma~\ref{lem:endpoint}
then preserves the entire integrability equation for
both trajectories. Proposition~\ref{prop:slice}
reconstructs actual HYM flows on the same $d_{x_b}$.

For initial data in $h^{17/8}$, the common parabolic
realization gives classical equations at positive
times and continuity in that space at time zero.
The reconstructed metric equation and its initial
value follow by this continuation. Smooth initial
metrics are included.
Lemma~\ref{lem:relative}, applied on the larger
coefficient ball, gives \eqref{eq:relative-log}
and \eqref{eq:order-main}. The remaining estimates
are those of Proposition~\ref{prop:phase}.
Its $C^1$ statement combines with the smooth change
$v_0=u_0-\Psi(p,x_0)$ to prove part~(iii).
Its uniqueness class and uniformity give the final
assertions. Proposition~\ref{prop:metric-charts}
below supplies the stated interpretation in initial
metric charts.
\end{proof}

\subsection{Coverage of nearby metrics}
\begin{proposition}\label{prop:metric-charts}
There is $r_{\mathrm{in}}\in(0,r/2)$ such that for
$|x_b|<r_{\mathrm{in}}$ every fixed-determinant
metric sufficiently close to $h_p$ in $h^{17/8}$
has a unique local representation
\begin{equation}\label{eq:metric-chart}
 H^{\mathrm{in}}=(e^{a_0})^*(h_pe^{u_0}),\qquad
 a_0\in\ha,\quad u_0\in Q_ph^{17/8},
\end{equation}
with $a_0,u_0$ small,
$x_0=e^{a_0}x_b\in B_r$, and
$\|u_0-\Psi(p,x_0)\|_{h^{17/8}}<\delta$.
The representation and its inverse are smooth with
uniform local derivative bounds in $p$, independent
of the stabilizer of $x_b$.

The same conclusion holds in a relative chart about
a possibly large gauge $\bar g$ provided
$|\bar g x_b|<r_{\mathrm{in}}$.
The chart uses the metric pulled back by $\bar g^{-1}$
and has constants independent of the absolute norm
of $\bar g$.
\end{proposition}

\begin{proof}
In an $h_p$-unitary presentation consider
\[
 \mathcal J_p(a,u)=\log(e^a e^u e^a):
 \ha\oplus Q_ph^{17/8}\longrightarrow h^{17/8}(\Herm_0E_0).
\]
The matrix product is positive Hermitian with
determinant one. Smooth composition in the H\"older
algebra makes $\mathcal J_p$ smooth. Its derivative
at zero is $(a,u)\mapsto2a+u$, with uniformly bounded
inverse $z\mapsto(\frac12P_pz,Q_pz)$.
The parameter-dependent inverse-function theorem
gives a common metric neighborhood.

The representation action near the identity sends
$B_{r_{\mathrm{in}}}$ into $B_{2r_{\mathrm{in}}}$.
Choose $r_{\mathrm{in}}$ after the tracking threshold
so that $C(2r_{\mathrm{in}})^2<\delta/2$ in the graph
bound. Reduce the metric neighborhood to ensure
$\|u_0\|<\delta/2$. The triangle inequality gives
the required initial normal displacement.
For a large gauge, apply the same factorization to
the pulled-back metric at coefficient $\bar g x_b$,
and return to the original bundle with
$g_0=e^{a_0}\bar g$. The estimates are in the relative
frame and introduce no factor $\|\bar g\|$.
\end{proof}

For integrable $x_b$, restriction of the ambient
phase map to these metric charts is $C^1$.
Confinement remains the hypothesis needed for its
all-time geometric interpretation.

\subsection{Special cases and scope}
For diagonal multiplicity groups, the slice equations
reduce to those used in the two-step metric response
and four-line phase constructions
\cite[Section~6]{wuebben3},
\cite[Sections~3--5]{wuebben8}.
In a positive diagonal frame $g=e^Y$, right group
differences become differences of $Y$.
Theorem~\ref{thm:main} therefore recovers their basic
confined phase and relative metric estimates when
the same graph construction is chosen. Their sharper
size factors and differentiated refinements remain
separate estimates.

The confined wall and triangular unstable trajectories
of \cite[Sections~3--5]{wuebben9} also fit the present
regime after the stated radius reductions. Their
gauges may diverge, while their harmonic coefficients
stay small. The relative reconstruction lemma explains
why this does not obstruct metric tracking in the
general noncommuting setting. Acyclic-support coercivity,
global entry and confinement, and precise logarithmic
or linear drift require the additional arguments of
those special families.

The stationary balancing and harmonic quotient comparisons
of \cite{wuebben14,wuebben15} do not supply a parabolically
invariant graph. The present proof uses the full Kuranishi
construction and moving stable-factor geometry, but
constructs its dynamical graph and phase independently
of those stationary comparisons. No assertion about
an evolving moduli-space metric follows merely from
the relative metric estimate proved here.
